\documentclass[11pt,a4paper]{article}

% Simple academic layout, following the original proposal.
\usepackage[T1]{fontenc}
\usepackage[utf8]{inputenc}
\usepackage[margin=1in]{geometry}
\usepackage{lmodern}
\usepackage{microtype}
\usepackage{amsmath,amssymb}
\usepackage{enumitem}
\usepackage[hidelinks]{hyperref}
\setlist[itemize]{topsep=5pt,itemsep=3pt,parsep=0pt}
\setlist[enumerate]{topsep=5pt,itemsep=4pt,parsep=0pt}
\setlength{\emergencystretch}{2em}
\clubpenalty=10000
\widowpenalty=10000

\title{Modifications to PointNeuron\\[0.4em]
\large Revised Research Proposal Following the PointNeuron1.0 Evaluation}
\author{}
\date{October 8, 2026}

\begin{document}
\maketitle

\noindent\textbf{Scope.}
This proposal retains the architectural directions of the original document and refines them using the completed PointNeuron1.0 diagnostic analysis. The immediate next step is external evaluation of the frozen baseline. Subsequent modifications are proposed experiments whose implementation and priority remain subject to discussion with Dr.~Arij.

\section{Semi-Automated Voxel-to-Point Transformation}

\subsection{Motivation}
Converting a neuron volume into a point cloud must preserve thin processes, terminals and branch junctions while limiting noise and computational cost. The PointNeuron1.0 diagnosis shows that input coverage interacts with alignment, proposal selection and graph construction. Samples 3, 4 and 6 contain out-of-bounds annotated nodes and trigger the alignment guard; these data/alignment problems must be distinguished from foreground thresholding failures~\cite{report,diagnostics}.

We propose investigating whether an image-derived reference skeleton can guide sampling toward informative neuronal structure. An automatically extracted skeleton may itself omit thin processes or introduce incorrect connections, so its benefit must be evaluated.

\subsection{Proposed Approach}
\begin{enumerate}
    \item \textbf{Reference skeleton extraction:} Apply a classical method, such as TEASAR or Kimimaro, to a segmented volume. Treat the extracted skeleton as an image-derived reference or pseudo-label, while retaining manually annotated SWC as the evaluation reference.
    \item \textbf{Geometry-guided sampling:} Compute each foreground point's distance to the reference skeleton and test whether weighted sampling preserves relevant structure better than the current protocol under the same point budget.
\end{enumerate}
One candidate weighting rule is
\begin{equation}
    w_i = \exp(-\alpha d_i), \qquad
    p_i = \frac{w_i}{\sum_j w_j},
\end{equation}
where $d_i$ is distance to the reference skeleton and $\alpha>0$ controls sampling concentration. This is a proposed rule to evaluate, rather than a setting selected for PointNeuron1.0.

\subsection{Implementation Details}
\begin{itemize}
    \item \textbf{Preprocessing:} Verify volume/SWC alignment, voxel spacing and intensity conventions. Record foreground thresholds, voxel counts and cap failures separately from alignment failures.
    \item \textbf{Patching:} Each Gold166 sample is a whole neuron volume. Local spatial point-cloud patches are formed within that neuron, and predictions are aggregated into whole-volume coordinates before graph reconstruction. The current cubic crop has radius 96, giving an interior width of 193 voxels at integer centers. Inference uses stride 96, at least 256 foreground voxels per patch and at most 256 coverage-selected patches. Completed promoted runs contained 3--220 patches~\cite{report}.
    \item \textbf{Sampling comparison:} Hold downstream components fixed. Compare terminal, branch and segment coverage, runtime and memory under matched point budgets.
    \item \textbf{Learnable refinement:} If a differentiable sampler is tested, define its learned parameters, gradient estimator and inference rule. A Chamfer-style loss against the reference may be explored, accounting for errors inherited from that reference.
    \item \textbf{Annotation use:} Annotated SWC may supervise training and evaluate predictions, but must not guide sampling on held-out test neurons. Manual correction requires a separate assisted protocol and a recorded intervention budget.
\end{itemize}
Training patches contain 512 points, whereas current full-volume inference aggregation defaults to 2048 sampled input points per patch; the latter was not serialized in saved inference archives. This difference could be tested in a future controlled ablation if unexplained failures remain, without attributing any observed failure to it.

\subsection{Expected Benefits}
The intended benefits are better retention of thin processes and junctions, fewer redundant input points and improved downstream coverage within a fixed computational budget. Sampling near an imperfect reference could instead reinforce missing branches or false connections.

\subsection{Discussion}
Reference-guided sampling combines classical image analysis with learned point-cloud processing. Its value depends on reference quality and whether improved input coverage survives proposal selection and graph reconstruction.

\section{Enhanced Point Cloud Encoder}

\subsection{Motivation and Overview}
The original framework uses DGCNN/EdgeConv. We propose evaluating PointNeXt as an alternative for extracting multiscale geometric features. The objective is to determine whether it better represents thin branches, terminals and junctions in this reconstruction task. Accuracy and efficiency gains must be measured rather than assumed from other point-cloud applications.

\subsection{Proposed Approach}
\begin{itemize}
    \item \textbf{Backbone replacement:} Compare PointNeXt with the existing encoder while retaining the sampling protocol, proposal head and graph procedure wherever interfaces permit.
    \item \textbf{Multiscale features:} Evaluate proposal coverage across process thicknesses and spatial densities.
    \item \textbf{Controlled comparison:} Use the same neuron-level splits and a documented training budget. Record necessary interface and optimization changes.
\end{itemize}

\subsection{Optional Transformer Fusion Module}
A lightweight attention module could subsequently supplement local features with longer-range context within a patch. Introduce it after the encoder comparison so its contribution and resource cost can be measured independently. Whole-neuron consistency will still depend on reconciling predictions from overlapping patches.

\subsection{Integration and Implementation Details}
The encoder produces features for the proposal module. Comparisons should report proposal, terminal and branch coverage, final geometry and topology, parameter count, runtime and peak memory. Neural feature extraction supports backpropagation; discrete sampling and graph operations require separate treatment.

\subsection{Expected Benefits}
The hypothesis is that improved multiscale features could increase node localization accuracy and coverage. A useful improvement must persist in the exported reconstruction and justify its computational cost.

\subsection{Discussion}
The initial experiment isolates the encoder. Optional transformer fusion and combined architectural changes follow only if this comparison identifies a measurable benefit.

\section{Advanced Skeleton Proposal Module}

\subsection{Motivation and Overview}
Node placement and selection affect connectivity differently across samples. Sample 22 had learned-node edge F1 approximately 0.110 versus oracle-node F1 approximately 0.890, with zero oracle bridges. In sample 19, segment coverage fell from 0.986 to 0.409 after selection~\cite{report,diagnostics}. These observations motivate separate investigations of selection and proposal generation.

\subsection{Proposed Approach}
\begin{enumerate}
    \item \textbf{Selection analysis:} Compare selectors using the same saved proposals and matched node budgets. Track coverage loss, skeleton precision, spatial distribution and downstream graph quality.
    \item \textbf{DETR-inspired proposals:} Evaluate a transformer decoder with learned queries predicting a set of coordinates, objectness scores and radii~\cite{detr}.
\end{enumerate}

\subsection{Implementation Details}
\begin{itemize}
    \item \textbf{Targets:} Define annotated endpoints and branch points, centerline samples at a specified spacing, or a documented combination. Raw SWC sampling density should not implicitly define the detection task.
    \item \textbf{Query budget:} Specify predictions per patch, unmatched-query supervision and treatment of target sets exceeding the budget.
    \item \textbf{Matching:} Define a bipartite matching cost and coordinate, objectness and radius losses, including units and handling of unavailable radius annotations.
    \item \textbf{Aggregation:} Reconcile overlapping-patch predictions in whole-volume coordinates. Set prediction within a patch does not establish uniqueness across patches.
\end{itemize}
A candidate objective is
\begin{equation}
    \mathcal{L}_{\mathrm{proposal}}
    = \lambda_{\mathrm{obj}}\mathcal{L}_{\mathrm{obj}}
    + \lambda_{\mathrm{pos}}\mathcal{L}_{\mathrm{pos}}
    + \lambda_{\mathrm{rad}}\mathcal{L}_{\mathrm{rad}},
\end{equation}
with regression evaluated on matched targets and objectness supervision defined for matched and unmatched queries. Loss choices and weights would be selected on validation data.

\subsection{Expected Outcomes and Benefits}
The intended outcomes are improved localization and endpoint/branch coverage with fewer redundant proposals. DETR avoids conventional NMS in its original formulation; this adaptation must test whether suppression can be reduced. Whole-neuron deduplication may remain necessary.

\subsection{Discussion}
Better candidates provide little benefit if selection discards them. A fixed query budget could also reproduce the existing coverage bottleneck. Evaluate retained coverage and final reconstruction quality, rather than training loss alone.

\section{Improved Connectivity Prediction}

\subsection{Motivation and Overview}
Accurate nodes are insufficient when candidate connections lack support or graph assembly introduces incorrect edges. Oracle-node reconstructions of samples 7, 16, 23 and 24 required 124, 104, 73 and 82 bridges, respectively. These cases motivate examining foreground support and graph construction together; they do not establish either as the sole cause. The existing connectivity GAE underperformed the initializer, so a new learned module requires a controlled comparison~\cite{report,diagnostics}.

\subsection{Proposed Approach}
\begin{itemize}
    \item \textbf{Candidate graph:} Document candidate construction and measure candidate-edge recall before learning scores. A scorer cannot recover an excluded target edge.
    \item \textbf{GATv2 comparison:} Compare a simple edge scorer with a GATv2 alternative using the same nodes, candidate edges and targets.
    \item \textbf{Evidence features:} Investigate endpoint confidence, Euclidean displacement, radius relations, directional continuation, snapping distance and foreground-geodesic support.
    \item \textbf{Assembly:} Evaluate edge predictions before tree construction and reconstruction quality after export.
\end{itemize}

\subsection{Implementation Details}
The current target is induced by mapping graph nodes to annotated SWC and deriving connectivity from SWC path distances. It can supervise training and evaluate predictions; GT-derived nodes or candidate sets remain labeled oracle diagnostics and are not deployed inference inputs.

Represent missing or unreachable geodesic support explicitly instead of treating a fallback cost as an observed path length. Topology-aware penalties may be tested alongside edge classification, with each penalty precisely defined and related to the decoding procedure.

\subsection{Bridge Count and Edge Validation}
A \textbf{bridge} is a fallback edge used to connect components when constrained foreground-geodesic construction cannot naturally produce a single connected tree. It is not a biological neuron branch. Bridges may be necessary for a valid SWC tree, but high counts indicate dependence on fallback connections. Interpret counts jointly with bridge hit rate and reachable-edge fraction.

Distinguish three forms of evidence:
\begin{enumerate}
    \item \textbf{Annotation agreement:} Matching the GT-induced connectivity derived from annotated SWC.
    \item \textbf{Image support:} Foreground paths and geometric continuation supporting a connection.
    \item \textbf{Independent biological assessment:} Expert review or other independent anatomical evidence supporting a reconstruction.
\end{enumerate}
Annotation agreement and image support are useful evidence, but neither independently proves biological plausibility. A valid tree can still contain incorrect connections.

\subsection{Expected Outcomes and Benefits}
The intended outcome is improved edge precision and recall with fewer unsupported bridges. GATv2 and topology penalties are candidate mechanisms rather than guarantees of improvement or biological correctness.

\subsection{Discussion}
Any gain must survive graph assembly. Evaluate the trade-off between forcing connectivity and retaining disconnected but better-supported components, documenting any fallback edges required for SWC export.

\section{End-to-End Integration and Training Strategy}

\subsection{Overview and Motivation}
The proposed pipeline retains the original sequence: voxel-to-point transformation, encoding, skeleton proposals, connectivity prediction and SWC export. Integration follows individual comparisons so improvements can be attributed to specific changes.

\subsection{Pipeline Integration}
\begin{enumerate}
    \item Extract foreground points and form local patches within each neuron volume.
    \item Encode patches and predict local skeleton candidates.
    \item Aggregate candidates into whole-volume coordinates and select graph nodes.
    \item Construct candidate edges, estimate connectivity and assemble the reconstruction.
    \item Export SWC and retain diagnostics from each stage.
\end{enumerate}

\subsection{Loss Functions and Optimization}
Where compatible, a future objective could combine sampling, proposal and connectivity losses:
\begin{equation}
    \mathcal{L}
    = \lambda_s\mathcal{L}_{\mathrm{sampling}}
    + \lambda_p\mathcal{L}_{\mathrm{proposal}}
    + \lambda_e\mathcal{L}_{\mathrm{edge}}
    + \lambda_t\mathcal{L}_{\mathrm{topology}}.
\end{equation}
This describes a possible objective, not proof of complete pipeline differentiability. Hard thresholding, discrete selection, candidate construction and export require explicit treatment.

\subsection{Training Strategy}
\begin{itemize}
    \item \textbf{Staged comparisons:} Begin with frozen inputs for selection and connectivity experiments. Train an encoder or proposal alternative when supported by diagnostic evidence.
    \item \textbf{Ablations:} Test patch point count, sampling, encoder, proposal head, selection and connectivity separately before combining changes.
    \item \textbf{Validation:} Choose optimization settings and loss weights on validation neurons. Document augmentation and preserve consistency between points, radii and annotations.
    \item \textbf{Joint fine-tuning:} Consider it after individual benefits are established and gradient paths through participating modules are defined.
\end{itemize}

\subsection{Expected Outcomes and Benefits}
Staged integration should provide interpretable evidence about which modules improve reconstruction. Joint optimization may subsequently offer additional benefits that must be measured.

\subsection{Discussion}
Training loss, architectural complexity and removal of post-processing are not sufficient evidence of superior reconstruction. Acceptance should depend on held-out geometry, connectivity, bridge dependence and resource use.

\section{Evaluation and Future Work}

\subsection{Evaluation Metrics}
\begin{itemize}
    \item \textbf{Data and preprocessing:} Alignment status, foreground thresholds and counts, cap satisfaction and execution failures.
    \item \textbf{Proposals and selection:} Candidate and selected counts, skeleton precision, terminal/branch coverage and coverage loss after selection.
    \item \textbf{Connectivity:} GT-induced edge precision, recall and F1; candidate reachability; bridge count and hit rate.
    \item \textbf{Geometry:} Exported-SWC point-distance precision, recall and F1, plus available approximate ESA/DSA/PDS metrics with their implemented definitions, distinguished from official benchmark metrics.
    \item \textbf{Validity and resources:} SWC validity, root count, runtime, peak memory and manual intervention.
\end{itemize}
Report topology F1 separately from spatial F1. Document thresholds, sampling spacing, units, denominators and missing values.

\subsection{Datasets}
Gold166 provides the existing PointNeuron1.0 evidence. On the 63 evaluated stable-domain samples (100--162), the promoted baseline achieved mean edge F1 0.8295, mean bridge count 1.444 and local point-distance SWC F1 0.8150. These are local results rather than external-test results or direct comparisons with the original paper benchmark~\cite{report}.

The proposed next dataset is Allen Cell Types, conditional on compatible paired microscopy stacks and SWC reconstructions~\cite{allen}. Verify stack completeness, spacing, coordinate alignment, intensity polarity and annotation scope. A projection or SWC alone is insufficient for the current volume input. The initial external evaluation uses frozen weights and reconstruction settings; any later adaptation is reported separately.

\subsection{Experimental Setup}
\begin{itemize}
    \item \textbf{Splits:} Keep patches from the same neuron in one split and group related specimens or donors where applicable.
    \item \textbf{Comparators:} Match inputs, candidate pools, node budgets and computational budgets where relevant.
    \item \textbf{Decision rules:} Agree on primary outcomes and acceptable trade-offs before examining held-out results.
    \item \textbf{Reporting:} Retain every requested sample, including exclusions and execution failures. Carry dataset identity, species, brain region and cortical layer when available; leave unrecorded metadata null.
\end{itemize}

\subsection{Results and Failure Analysis}
Future results should include per-sample evidence alongside aggregate metrics. The diagnostic schema distinguishes execution status from \texttt{CLEAR}/\texttt{REVIEW}, which indicate diagnostic review needs. Preserve separate alignment, foreground preprocessing, proposal, selection and graph/geometry evidence, with conservative attribution for ambiguous cases~\cite{diagnostics}.

Inspect representative failures using input overlays, candidate-to-selected coverage, edge support and final SWC geometry. Explain which failure mechanisms a modification addresses and where it remains insufficient.

\subsection{Future Work}
\begin{enumerate}
    \item Verify Allen compatibility and define the external-test protocol.
    \item Evaluate frozen PointNeuron1.0 using the same diagnostic schema.
    \item Choose one intervention motivated by observed failures and compare it against a controlled baseline.
    \item Combine supported changes and repeat held-out evaluation, adding independent morphology review if biological plausibility is to be claimed.
\end{enumerate}
Transformer fusion and broader architectural integration remain later options contingent on these comparisons.

\subsection{Discussion}
External validation and controlled ablations should establish transfer performance and explain reconstruction errors. Their findings will determine which elements of the architectural proposal warrant further development.

\section{Conclusion}
This revised proposal preserves the original directions of adaptive sampling, PointNeXt encoding, DETR-inspired proposals and GATv2 connectivity while grounding them in the completed PointNeuron1.0 diagnosis. The proposed contribution is a validated account of how input coverage, selection and edge construction affect reconstruction, followed by targeted improvements supported by controlled experiments.

The immediate decisions for Dr.~Arij are the scope of external validation, the first intervention to prioritize and the evidence needed to accept an improvement. PointNeuron2.0 remains a proposed research direction whose implementation follows that review.

\begin{thebibliography}{9}
\bibitem{report}
PointNeuron project. \emph{PointNeuron1.0 Reimplementation and Evaluation}.
Completed diagnostic closure and schema cleanup.
Local report: \path{docs/PointNeuron1_Report.pdf}.

\bibitem{diagnostics}
PointNeuron project. \emph{Per-sample diagnostics and source manifest}.
\path{tmp/diagnostics/pointneuron1_sample_diagnostics.json};
\path{tmp/diagnostics/pointneuron1_sources.json}.
Schema: \path{docs/sample_diagnostics.md}.

\bibitem{detr}
N.~Carion et al. \emph{End-to-End Object Detection with Transformers}.
European Conference on Computer Vision, 2020.
\url{https://www.ecva.net/papers/eccv_2020/papers_ECCV/html/832_ECCV_2020_paper.php}.

\bibitem{allen}
Allen Institute. \emph{Cell Types Database: API}.
Official documentation on morphology and image access.
\url{https://community.brain-map.org/t/cell-types-database-api/3016}.
\end{thebibliography}

\end{document}
